\section{Conclusion}
\label{sec:conclusion}

This paper presented \textsc{Nexus}, a user-space KV-cache Memory Management Unit (MMU) designed to \emph{bypass} the tool schema prefill wall (for $P \le 256$; with a text-prefill fallback beyond) in agentic serving systems. By decoupling tool schemas from online attention computation, \textsc{Nexus} compiles schemas offline into binary Aeon Tool Blocks (\texttt{.atb}) and splices them directly into active mid-context positions. The system coordinates L1 Semantic Lookaside Buffer embedding scans, Left-Child Right-Sibling radix trie FSM constraint masking, and stride-aware memory mapping (the \texttt{.atb} block residency is zero-copy; the splice into the live sequence is a \texttt{seq\_cp} plus a fused tail recompute).

Measurements on Apple Silicon (Darwin arm64) running Qwen2.5-14B-Instruct show \textsc{Nexus} achieves a gateway routing+splice latency P50 of $\approx 160\,\text{ms}$ (regenerated 2026-06-20; this excludes first-token generation), and an end-to-end TTFT $\approx 2.34\times$ faster than retrieve-and-prefill measured like-for-like, at $89.0\%$ tool-routing accuracy (a $\approx 3$-point trade-off vs.\ $B_{RP}$'s 0.92). Critical-path operations in the C++ hot path execute on contiguous, cache-aligned, pre-reserved arenas with no dynamic heap allocation. \emph{Erratum:} the previously reported $237.4\,\text{ms}$ TTFT / $6.3\times$ figure did not reproduce; see \texttt{docs/AUDIT\_2026-06-20\_LIVE.md}.

Our investigation of deep-context splicing boundaries and the causal isolation constraint defines the mathematical and physical limits of offline KV splicing under RoPE phase drift. \textsc{Nexus} establishes a predictable, high-performance systems foundation for real-time, multi-domain conversational agents.
